\subsection{Kenmerken voor de detectie van phishingmails}
\label{subsec:Kenmerken voor de detectie van phishingmails}
Phishing is een vorm van cybercriminaliteit waarbij een aanvaller het vertrouwen van het slachtoffer probeert te winnen door zich voor te doen als een betrouwbare partij. Op die manier tracht de aanvaller het slachtoffer ertoe te bewegen gevoelige informatie prijs te geven of toegang te verlenen tot beschermde gegevens. In tegenstelling tot veel andere vormen van cyberaanvallen, die zich rechtstreeks op systemen of technische kwetsbaarheden richten, is phishing voornamelijk gericht op de menselijke factor en maakt het gebruik van sociale manipulatie om slachtoffers tot bepaalde handelingen aan te zetten \autocite{Hong2012}.

De automatische detectie van frauduleuze e-mails vertrekt vanuit de vraag welke kenmerken een phishingmail onderscheidt van een legitieme. Dit hoofdstuk behandelt net die kenmerken die reeds werden onderzocht in voorgaand onderzoek.

\subsubsection{Overzicht en classificatie van detectiekenmerken}
\label{subsubsec:Overzicht en classificatie van detectiekenmerken}
Automatische phishingdetectie in e-mails is de voorbije jaren sterk geëvolueerd. Vroege detectiesystemen steunden op handmatig gedefinieerde kenmerken die als invoer dienden voor machinelearningmodellen. Zo introduceerden \textcite{Fette2007} PILFER, een classificatiemodel dat op basis van tien handmatig gekozen kenmerken bepaalt  of een e-mail phishing is. \textcite{Bergholz2010} bouwden hierop voort en combineerden  inhoudelijke en structurele kenmerken, zoals statistische modellen voor de onderwerpen van een e-mail en de sequentiële analyse van tekst en externe links. Later onderzoek maakt  daarnaast gebruik van machine- en deep-learningmodellen \autocite{Salloum2022}.